% =====================================================================
%  General Conclusion
% =====================================================================
\chapter*{General Conclusion}
\addcontentsline{toc}{chapter}{General Conclusion}
\markboth{GENERAL CONCLUSION}{GENERAL CONCLUSION}

The work documented in this thesis has covered the full CRISP-DM cycle of an industrial data engineering and machine learning project conducted at \hostorg, on a multi-tenant telematics fleet of approximately six hundred GPS devices spread across four customer organizations. The starting point was an operational ecosystem composed of three MySQL databases --- a raw hexadecimal telemetry stream populated by SIM-driven device frames, a decoded archive of \texttt{arch\_xxxxx} tables retaining about six months of readable data, and a \texttt{user\_xxxxx} master-data database carrying customers, vehicles and drivers --- on top of which the company already exposed a set of descriptive dashboards but no advanced analytical layer. The endpoint is a production-grade analytical platform deployed on a Microsoft Azure virtual machine, refreshed every five minutes, equipped with an explicit MySQL-to-PostgreSQL conversion bridge, a cleaning rule engine, a dimensional warehouse, a registry of thirty-five behavioural features, a per-tenant K-Means segmentation, an Isolation Forest risk score, an API and a business intelligence dashboard.

The first half of the project has invested heavily in the data engineering layer. A dedicated bridge converts the \texttt{.sql} dump files received from \hostorg{} --- covering the slices of the \texttt{arch\_xxxxx} archive and the \texttt{user\_xxxxx} master data relevant to the analytical scope --- into a PostgreSQL-compatible representation. Eleven data quality issues were identified during the data understanding phase and were formalized as cleaning rules executed at the warehouse boundary. A dimensional model with five dimensions, a bridge and eleven facts has stabilized the silver layer, on top of which the analytical marts and the unified ML view have been built. The watermark-driven incremental pipeline, hosted on the Azure VM and orchestrated by Prefect, has reduced the per-cycle processing time to a few seconds while preserving idempotency and replayability, and has done so without ever touching the operational MySQL workload of the company.

The second half of the project has built the analytical and serving layers. After a comparative study covering rule-based, statistical, classical machine-learning and deep-learning candidates, a pair K-Means + Isolation Forest, fitted per tenant, has been retained. The behavioural profiles obtained align with the operational vocabulary of the fleet managers; the rescaled risk score and its three operational bands feed the business intelligence dashboard and the API consumed by the consumer applications. A layered monitoring covers the pipeline health, the data quality and the model drift, with alerting hooked to the operations channel. Together, these deliverables move the analytical centre of gravity of \hostorg{} from a purely descriptive surface --- \emph{what happened?} --- to a forward-looking decision-support layer --- \emph{which devices are likely to misbehave next month?}

A particular care has been given to the secure remote collaboration on the Azure infrastructure, dictated by the binomial nature of the project. Both members of the binomial worked from separate personal workstations and shared the analytical environment through their respective SSH key pairs. The PostgreSQL ports are not exposed to the public internet; access is mediated by SSH tunnels, themselves protected by a strict IP allowlist applied at the Azure Network Security Group level and limited to the binomial's workstations and the office network of \hostorg. Per-member personal accounts coexist with a dedicated automation account on the VM, under controlled file permissions; the database service account is provisioned with the minimal privileges required by the pipeline.

Three lessons emerge from the project. \emph{First}, the early formalization of the cleaning rules has been the single most effective decision: it has prevented the propagation of source-side anomalies into the modeling layer and has made every transformation auditable, an essential property when the source is an external dump file rather than a directly accessible database. \emph{Second}, the per-tenant fitting protocol has been a decisive factor in obtaining operationally meaningful clusters; a global model would have been dominated by the most active tenant and would have failed to capture the heterogeneity of the regimes anticipated in the data understanding phase. \emph{Third}, the strict separation between the operational MySQL stack of \hostorg{} and the analytical PostgreSQL warehouse of the project has been more than an administrative convenience: it has freed the analytical work to iterate on cleaning rules, feature definitions and modeling choices without ever putting the customer-facing platform of the company at risk.

Three perspectives are open for future iterations. The first one is the integration of a labelled incident feed from the operations team of \hostorg, which would unlock a supervised classification track and a quantitative calibration of the risk score against ground truth. The second one is the migration of the dump-based ingestion towards a streaming back-end (Kafka, Redpanda or a managed equivalent), in order to bring the latency below the minute and to support event-driven downstream consumers; this would naturally complement the existing streaming-style ingestion of the raw hexadecimal telemetry on the operational side. The third one is the introduction of a hierarchical or multi-task model capable of leveraging the structural similarities across tenants while preserving the tenant-specific regimes already identified.

Beyond the technical deliverables, the project has been a formative experience on the trade-offs between elegance and operational sustainability, and between methodological rigour and the very concrete constraints of a binomial collaboration. Building a pipeline that runs successfully once is comparatively easy; building a pipeline that runs every five minutes for several months without manual intervention, on an infrastructure jointly operated by two student engineers from separate workstations, requires a different kind of engineering --- one in which observability, idempotency, explicit failure modes and disciplined access control carry as much weight as algorithmic sophistication. The platform delivered in this thesis aims to incarnate that discipline.

\cleardoublepage
